\chapter{Getting started}\label{ch:start}

\section{Opening PlanWise}

PlanWise is a web application. It consists of a backend (the scheduling service and its database) and the
web interface. Ask your administrator for the address. On a local installation it is:

\centerline{\code{http://localhost:3000}}

Any current browser works (Chrome, Edge, Firefox); a screen at least 1280 pixels wide is recommended,
because the Gantt chart uses the full width.

\begin{note}
For a local installation, an engineer starts both parts with one command from the project folder:
\code{powershell -ExecutionPolicy Bypass -File scripts/dev.ps1}. With Docker, \code{docker compose up
--build} does the same. The API documentation is then available at
\code{http://localhost:8000/api/v1/docs}.
\end{note}

\subsection{Access and roles}

In a development installation no sign-in is needed. In production the backend is protected by keys:

\begin{itemize}
  \item an \textbf{operator} key allows every action (editing data, generating and committing schedules);
  \item a \textbf{reader} key allows viewing only. A reader sees every page, but actions that change data are
        refused, and the schedule editor reports ``This account is read-only and cannot edit schedules.''
\end{itemize}

Every change is recorded with the actor that made it, so committed decisions, set-up versions and imports
can always be traced back.

\section{The screen layout}\label{sec:layout}

Every page shares the same frame (Figure~\ref{fig:today} in the next chapter shows it with callouts).

\begin{description}[style=nextline, leftmargin=1.2em, font=\bfseries\color{brand}]
  \item[Sidebar (left)] The main navigation, grouped as \ui{Schedule} (Today, Planning periods, Order
        book), \ui{Plant set-up} (Factory) and \ui{Advanced} (JSON scenarios, Runs, Aggregate planning).
        The Advanced group is folded by default because planners rarely need it; click its title to open it.
        Your choice is remembered in this browser.
  \item[New button] A shortcut menu for the most common creations (Figure~\ref{fig:new_menu}).
  \item[Factory name] Opens the Factory page. When several plants exist it lets you switch plant.
  \item[Quick search] Opens the command palette (\kbd{Ctrl}\,\kbd{K}) to jump to any page, planning
        period, scenario, recent run or action (Section~\ref{sec:palette}).
  \item[Status box (bottom left)] \ui{System online} when the backend answers, the number of solver runs in
        progress (or \ui{Queue is empty}), the database response time and the time of the latest schedule.
        It turns to \ui{API unreachable} when the backend cannot be reached.
  \item[Top bar] The breadcrumb of the current page, the environment badge (for example
        \ui{development}), the language switch (\ui{EN}/\ui{VI}) and the light/dark theme menu.
\end{description}

\shotf[0.36\linewidth]{new_menu}{The \ui{New} menu: create a planning period, import orders, add one order, or (Advanced) import a JSON scenario, start a single run or create an aggregate plan.}{new_menu}

\section{Language and theme}

The whole interface is bilingual. Click the language button (\ui{EN} or \ui{VI}) in the top bar and pick
\ui{English} or \ui{Tiếng Việt}. The choice is stored in a browser cookie, so it survives reloads, and it
also changes date formats and number formatting. Codes entered by users (order numbers, customer names,
machine names) are shown as they were entered.

The theme menu next to it offers \ui{Light}, \ui{Dark} and \ui{System} (follow the operating system). See
Section~\ref{sec:theme} for a dark-theme screenshot.

\section{The first time: setting up a plant}\label{sec:first-setup}

When no factory exists yet, Today, Order book and Planning periods all show a short message with a
\ui{Set up the factory} button. On the Factory page you can create the plant from a \term{schema-5 JSON
file}, for example the research dataset \code{dataset/wheel-factory-2weeks/raw\_input.json}:

\begin{steps}[Create the factory from a file]
\begin{enumerate}
  \item Open \ui{Factory} in the sidebar.
  \item Click \ui{Choose a schema-5 JSON file} and select the file.
  \item Check the proposed factory name and click \ui{Create the factory}.
\end{enumerate}
\end{steps}

The file is split by how often its parts change: plant data becomes factory set-up version~1, its orders go
to the order book, and a first planning period is created. From then on, each period only needs orders and
opening stock; the set-up is edited on the Factory page (Chapter~\ref{ch:factory}).

\begin{tip}
A file that is not a schema-5 factory document (it needs stages, a shift template, machines and products)
is rejected with the message ``Not a schema-5 factory file'' and nothing is created.
\end{tip}
